\chapter{Tools and Monitored VM Onboarding Inventory}

This appendix summarizes the tools/setup material inspected for the report. It is included as operational documentation because ARIA depends on monitored machines sending logs, alerts, runtime events, and metrics to the central stack.

\section{Central Environment Scripts}

\begin{longtable}{@{}p{0.30\textwidth}p{0.60\textwidth}@{}}
\caption{Central environment setup files}\\
\toprule
\textbf{File} & \textbf{Role} \\
\midrule
\endfirsthead
\toprule
\textbf{File} & \textbf{Role} \\
\midrule
\endhead
\sourcefile{tools/tools/setup_script.sh} & Wrapper that delegates to the maintained central setup flow. \\
\sourcefile{tools/tools/setup_script_telegraf.sh} & Central orchestration script using component flags for selected services. \\
\sourcefile{tools/tools/siem_setup.sh} & Elasticsearch, Kibana, Filebeat, TLS and system log forwarding setup. \\
\sourcefile{tools/tools/wazuh_setup.sh} & Wazuh Manager setup and Wazuh alert integration. \\
\sourcefile{tools/tools/suricata_setup.sh} & Suricata IDS setup and EVE JSON output configuration. \\
\sourcefile{tools/tools/setup-falco-server-elastic.sh} & Falco/Falcosidekick setup for runtime event forwarding. \\
\sourcefile{tools/tools/telegraf_setup.sh} & Telegraf metric collection and forwarding setup. \\
\sourcefile{tools/tools/detection_rules_setup.sh} & Detection-rule and visualization object setup. \\
\sourcefile{tools/tools/hardening_setup.sh} & Host hardening actions such as firewall and SSH-related controls. \\
\bottomrule
\end{longtable}

\section{Monitored VM Bootstrap Flow}

The monitored VM bootstrap script is located under \sourcefile{tools/Script that i should inject in VMs to monitor them/bootstrap_monitored_vm.sh}. A copy also exists under the backend scripts directory for project distribution. Its purpose is to prepare a VM that will be observed by ARIA.

\begin{enumerate}[leftmargin=1cm]
\item Confirm that the target is a monitored VM and not the central ARIA/monitoring server.
\item Collect the VM name, environment label, IP address, central Elasticsearch URL or IP, credentials, optional certificate path, and optional Wazuh manager configuration.
\item Test central Elasticsearch reachability before installing agents.
\item Install selected components: Wazuh Agent, Filebeat, Suricata, Falco/Falcosidekick, and Telegraf.
\item Write configuration files with monitored asset fields such as host name, asset name, environment, and role.
\item Start and verify services.
\item Query the central stack to check that expected index families are visible.
\item Print a JSON payload that can be used to create or edit the ARIA monitored asset.
\end{enumerate}

\section{Sensitive Values}

Screenshots and appendices must never expose passwords, API keys, private SSH values, or full secret headers. The report should use placeholders such as \codepath{<redacted>} when showing commands that require credentials.

\section{Validation Commands to Capture}

\begin{longtable}{@{}p{0.38\textwidth}p{0.52\textwidth}@{}}
\caption{Useful onboarding validation commands}\\
\toprule
\textbf{Command type} & \textbf{Purpose} \\
\midrule
\endfirsthead
\toprule
\textbf{Command type} & \textbf{Purpose} \\
\midrule
\endhead
\codepath{systemctl status wazuh-agent filebeat suricata falcosidekick telegraf --no-pager -l} & Prove that installed monitored-VM services are running or clearly identify which optional component is absent. \\
Elasticsearch health query with credentials hidden & Prove that the VM can reach the central stack. \\
Elasticsearch index listing for asset-specific patterns & Prove that logs, runtime events, and metrics are arriving. \\
ARIA asset source validation action & Prove that the backend can query source data for the asset. \\
ARIA alerts or metrics filtered by \codepath{asset_id} & Prove that telemetry is visible in the SOC workflow. \\
\bottomrule
\end{longtable}

\section{Known Setup Limitations}

\begin{itemize}[leftmargin=1cm]
\item The bootstrap script prints the ARIA asset payload but does not automatically register the asset in all environments.
\item Some components are optional; missing data for a source can be valid when that component was not installed or configured for the VM.
\item Central setup scripts are provisioning scripts and should be reviewed before use on an existing server.
\item TLS and certificate handling must be adapted to the deployment constraints and should not be bypassed for a production deployment.
\end{itemize}

